% SPDX-FileCopyrightText: Copyright (c) 2024-2026 Yegor Bugayenko
% SPDX-License-Identifier: MIT

\section{Principles}

\subsection{Preciseness Over Expressiveness}

We treat objects as the words of an OOP prose:
  the objects that reside at the top level of the root package
  and its sub-packages are the nouns of its vocabulary.
The more concise this vocabulary is and the less ambiguous,
  the cleaner the prose written with it.
Synonyms are especially harmful here,
  since two names for the same meaning make the prose more artistic
  but less strict, forcing the reader to guess which one carries the intent.
We therefore select global objects carefully,
  making sure that none of them duplicate each other,
  and that each delivers a clear and distinct meaning.

\subsection{Transparency Over Speed}

An atom is a function implemented by the runtime beneath \eolang{},
  by the virtual machine or the CPU itself,
  and therefore runs genuinely faster than its \eolang{} counterpart.
An atom is also opaque, though:
  neither the \eolang{} compiler nor any static analyzer
  can look inside it, reason about it, or optimize it.
A program saturated with atoms looks like a solid brick to these tools---
  it either works as it happens to work,
  or nothing can be done about it.
We instead want as many objects as possible to stay transparent to the compiler,
  which we expect to reason about them more powerfully
  than the programmers of atoms ever could.
Eventually, we expect an OO program written in \eolang{},
  resting on only a handful of atoms,
  to be optimized down to highly efficient assembly code.

\subsection{Reusability Over Complexity}

We treat our objects as the building blocks of larger OOP systems,
  and we want every block to be used and to be known to the \eolang{} compiler.
A programmer who reaches for them should never need to reinvent the wheel:
  we intend to supply everything,
  from array sorters and regular expression processors
  to file readers and writers.
The more reusable each object is, the weaker the temptation
  to implement something ``better'' elsewhere,
  and the fewer the private duplicates scattered across programs.
We accept that a few of our objects may grow complex,
  or carry functionality that feels redundant, and that is fine.
Later, once the \eolang{} compiler becomes powerful enough,
  it will know how to reduce that complexity and optimize it away at compile time.

\subsection{Decoration Over Composition}

Decoration is our primary method of extending functionality.
When an object lacks a feature,
  we expect the programmer to build a new object that looks like the original one,
  keeping the same interface but behaving differently.
Wrapping objects into composites with new interfaces and more powerful features,
  as the Adapter and Facade patterns suggest~\citep{gamma1994design},
  is another way to extend what already exists, and we discourage it:
  a new interface hides the original object from its users and from the compiler,
  while a decorator stays interchangeable with the object it decorates.

\subsection{Currying Over Polymorphism}

Suppose an object \ff{books} must tell whether it contains a given title.
The traditional object-oriented answer is polymorphism:
  \ff{books} exposes a \ff{contains} attribute
  and the caller writes \ff{books.contains "Foo"}.
Different kinds of \ff{books} implement it differently---
  an in-memory list scans its elements,
  while a database-backed one runs a \ff{SELECT} query.
The algorithm stays under the control of \ff{books},
  which is the strength of polymorphism,
  but every such feature enlarges the interface of \ff{books},
  and that interface keeps bloating as the code base grows.

The opposite answer moves \ff{contains} out of \ff{books}
  and makes it a free object that takes \ff{books} as an argument,
  so the caller writes \ff{contains books "Foo"}.
Now \ff{books} stays small,
  but polymorphism is lost:
  \ff{books} no longer decides how the search is performed.

We suggest a third answer that keeps both benefits.
The searching object is handed to the caller already bound to its collection,
  as a partially applied object,
  and the caller only completes it, as \ff{contains "Foo"}.
The interface of \ff{books} stays minimal,
  yet the concrete algorithm is still chosen by whoever assembled \ff{contains}.
This is currying, or partial application~\citep{barendregt2012},
  familiar from functional languages such as Haskell,
  brought into an object-oriented setting.
